\documentclass[10pt,a4paper]{amsart}
\usepackage[margin=19mm]{geometry}
\usepackage{amsmath,amssymb,mathtools}
\usepackage[scr=rsfs,cal=euler]{mathalfa}
\usepackage{xfrac}
\usepackage[unicode,hidelinks,bookmarks=false]{hyperref}
\newcommand{\ie}{i.e.\ }
\newcommand{\cf}{cf.\ }
\newcommand{\resp}{respectively\ }
\newcommand{\Subsection}{\subsection}
\providecommand{\nobreakdash}{\nobreak\hbox{-}}
\setlength{\emergencystretch}{2em}
\multlinegap0pt
\usepackage{amssymb}
\usepackage[shortlabels]{enumitem}
\usepackage{dsfont}
\newtheorem{lemma}{Lemma}
\def\gam{\gamma}
\title{Critical regularity and dissipativity for stochastic reaction-diffusion equations in Bochner spaces over spaces of continuous functions}
\author{Xuewei Ju, Xiaoting Tong}
\date{Source: arXiv:2604.13625v3 (CC BY 4.0)}
\begin{document}
\maketitle

\noindent\emph{Source and licence.} Critical regularity and dissipativity for stochastic reaction-diffusion equations in Bochner spaces over spaces of continuous functions. Version arXiv:2604.13625v3 (CC BY 4.0), distributed by arXiv under the Creative Commons Attribution 4.0 International licence, http://creativecommons.org/licenses/by/4.0/, as recorded in the arXiv licence metadata for this exact version and shown on the abstract page https://arxiv.org/abs/2604.13625v3. Full licence notice: \texttt{licenses/arxiv-2604.13625-CC-BY-4.0.txt}.\par\smallskip

\noindent\textit{Editorial scope.} Counted unit: the article's lemma le4.3 (source0.tex lines 1351--1364). The article's complete original proof of it is delivered with it (source0.tex lines 2154--2246, one \texttt{proof} environment of 93 lines, which the article places away from the statement). No mathematics is written, repaired or re-derived: the statement and the proof are the article's own lines, in the article's own notation, and no retained line is substituted or re-flowed.  No further source-local context is needed: every cross-reference of every retained line resolves inside the two delivered spans.  Nothing was omitted from any retained span, so the package carries no omission record.  No retained line contains a citation, so the wrapper carries no bibliography and no bibliography entry.  No retained line contains a figure environment, an image-inclusion command or drawing code, so no image is delivered and the package declares no asset dependency.  Editorial formatting only, in addition: an AMS \texttt{amsart} preamble that restates the article's own declarations and macros used by the retained lines, namely the environments \texttt{lemma} on separate counters, together with the standard packages those lines need.  The retained environments print in this document as Lemma 1 in order of appearance, whereas the article prints the same items with the numbering of its own full document.  The complete native source of the article as distributed in the arXiv e-print for this exact version is bundled unchanged as \texttt{sources/198488/source0.tex}.
\begin{lemma}\label{le4.3}
For any $d\in\mathbb N^+$, there exist parameters $\nu,\gamma\in(0,1/2)$ satisfying
$$
\begin{cases}
0<\frac{d}{2(q-1)}<1,\\[0.3em]
0<\frac{2\nu q}{q-2}<1,\\[0.3em]
0<\left(1+\frac{d}{2q}-\nu+\gam\right)\frac{q}{q-1}<1,\\[0.3em]
0<(\gamma+\frac{d}{2q})\frac{q}{q-1}<1,\\[0.3em]
q\gamma>1,\\[0.3em]
q\nu>\frac{d}{2},
\end{cases}
$$
if and only if $$q>d+6.$$ %Moreover, the lower bound $26$ is sharp.
\end{lemma}

\begin{proof}[Proof of Lemma \ref{le4.3}]
\textbf{Necessity.}
From the first inequality: $0<\frac{d}{2(q-1)}<1$ yields $q>1+\frac d2$.
From the second inequality: $0<\frac{2\nu q}{q-2}<1$. The left inequality is trivial for $\nu>0$. The right inequality gives
$$
\nu < \frac{q-2}{2q}=\frac12-\frac1q.
$$
Since $\nu\in(0,\frac12)$, we require $q>2$.
Condition $q\gamma>1$ implies $\gamma>\frac1q$. Combined with $\gamma<\frac12$, we need $q>2$.
Condition $q\nu>\frac d2$ implies $\nu>\frac{d}{2q}$.
Combined with $\nu<\frac12-\frac1q$, a necessary condition for existence of $\nu$ is
$$
\frac{d}{2q}<\frac12-\frac1q \implies d < q-2 \implies q>d+2.
$$
From $0<\big(\gamma+\frac{d}{2q}\big)\frac{q}{q-1}<1$, the left side is positive. The right-hand inequality:
$$
\gamma+\frac{d}{2q}<1-\frac1q \implies \gamma<1-\frac{d+2}{2q}.
$$
When $q>d+2$, one checks $1-\frac{d+2}{2q}>\frac12$, so the effective upper bound for $\gamma$ is $\frac12$.
Thus
$$
\gamma\in\Big(\frac1q,\,\frac12\Big).
$$
Next consider $0<\big(1+\frac{d}{2q}-\nu+\gamma\big)\frac{q}{q-1}<1$.
The left inequality is easily satisfied for large $q$. The critical constraint comes from the right inequality:
$$
1+\frac{d}{2q}-\nu+\gamma < 1-\frac1q
\;\Longrightarrow\;
\gamma -\nu < -\frac{d+2}{2q},
\quad \text{i.e.}\quad
\gamma < \nu -\frac{d+2}{2q}.
$$
Collect all hard constraints for $q>d+2$:
$$
\begin{cases}
\dfrac{d}{2q}<\nu<\dfrac12-\dfrac1q,\\[0.4em]
\dfrac1q<\gamma<\dfrac12,\\[0.4em]
\gamma < \nu -\dfrac{d+2}{2q}.
\end{cases}
$$
To have a pair $(\nu,\gamma)$, the interval for $\gamma$ must be non-empty:
$$
\frac1q < \nu -\frac{d+2}{2q}.
$$
This must hold for some admissible $\nu$. The maximal available $\nu$ is $\nu\nearrow \frac12-\frac1q$. Substitute this upper bound into the inequality to get a necessary condition:
$$
\frac1q < \left(\frac12-\frac1q\right)-\frac{d+2}{2q}.
$$
Multiply both sides by $2q>0$:
$$
2 < q -2 -(d+2) \implies q>d+6.
$$

\noindent\textbf{Sufficiency.}
Assume $q>d+6$. Choose
$$
\nu_0=\frac12-\frac1q-\varepsilon,\qquad
\gamma_0=\frac1q+\varepsilon,
$$
where $\varepsilon>0$ is chosen sufficiently small such that
$$
\varepsilon<\min\left\{\frac12-\frac1q,\ \frac{3}{2q},\ \frac14-\frac{d+4}{4q}\right\}.
$$
Such a choice is possible since $d\in\mathbb N^+$ gives $q>d+6\ge7$, hence $q>2$, and
$$
\frac14-\frac{d+4}{4q}>\frac14-\frac{d+4}{4(d+6)}=\frac{1}{4(d+6)}>0.
$$

We verify:
1. $\nu_0\in(0,1/2)$: since $\varepsilon$ is small, $\nu_0>0$ and $\nu_0<1/2$.
2. $\gamma_0\in(0,1/2)$: since $\varepsilon<1/2-1/q$, we have $\gamma_0<1/2$.
3. $q\gamma_0=1+q\varepsilon>1$.
4. $q\nu_0=(q/2)-1-q\varepsilon$. Since $q>d+6>d+2$, we have $q/2-1>d/2$, hence $q\nu_0>d/2$ for $\varepsilon$ small enough.
5. For the second inequality: $\frac{2\nu_0 q}{q-2}=\frac{q-2-2q\varepsilon}{q-2}<1$.
6. For the critical constraint, we require
$$
\gamma_0 < \nu_0-\frac{d+2}{2q}.
$$
Substituting $\nu_0,\gamma_0$ gives
$$
\frac1q+\varepsilon < \frac12-\frac1q-\varepsilon-\frac{d+2}{2q}
\iff
2\varepsilon < \frac12-\frac{d+4}{2q}.
$$
Since $q>d+6$, we have
$$
\frac12-\frac{d+4}{2q}>\frac12-\frac{d+4}{2(d+6)}=\frac{1}{d+6}>0.
$$
Our choice of $\varepsilon$ ensures this inequality holds.
7. The fourth inequality follows similarly.

Thus all conditions are satisfied for sufficiently small $\varepsilon>0$. Hence $q>d+6$ is sufficient.
\end{proof}

\end{document}
